% !PRINTED 332
\weberpage{332}
%<<CONT
2. \gesperrt{Nennen wir die Form $T$ die zu der Basis $(\alpha_\nu)$ von $\mathfrak{a}$ gehörige Form, so sind die zu verschiedenen positiven Basen desselben Ideals gehörigen Formen äquivalent.}

Bezeichnen wir die konjugierten Werte von $\alpha_s, \omega_s$ mit

\[
\alpha_{s,1},\ \alpha_{s,2},\ \dots,\ \alpha_{s,n};\quad \omega_{s,1},\ \omega_{s,2},\ \dots,\ \omega_{s,n}
\]

und setzen

\weq{10}{\alpha_{s,r} = \overset{\nu}{\Sigma}\, a_{s,\nu}\,\omega_{\nu,r},}

so ergibt sich

\weq{11}{(\Sigma \pm \alpha_{1,1}\,\alpha_{2,2}\,\dots\,\alpha_{n,n})^2 = N(\mathfrak{a})^2 \varDelta,}

wenn $\varDelta$ die Grundzahl des Körpers $\Omega$, nämlich das Determinantenquadrat:

\[
\varDelta = (\Sigma \pm \omega_{1,1}\,\omega_{2,2}\,\dots\,\omega_{n,n})^2
\]

ist (Bd.\,II, \S\,162).

Sind $\lambda_1, \lambda_2, \dots, \lambda_n$ die konjugierten Werte von $\lambda$, so ist nach (4)

\weq{12}{\lambda_r = \overset{\nu}{\Sigma}\, \alpha_{\nu,r}\, t_\nu}

eine lineare Substitution für die Variablen $t_\nu$ mit der Substitutionsdeterminante

\weq{13}{r = N(\mathfrak{a})\,\sqrt{\varDelta},}

und durch diese geht nach (5) die Form $T$ in das Produkt

\weq{14}{T' = [N(\mathfrak{a})]^{-1}\,\lambda_1\,\lambda_2\,\dots\,\lambda_n}

über, weil $N(\lambda)$ das Produkt der konjugierten Werte von $\lambda$ ist (Bd.\,II, \S\,151). Wir können hierauf die Invariantentheorie (Bd.\,I, \S\,65) anwenden, und wenn wir die \gesperrt{Hessesche} Determinante bilden:

\weq{15}{H = \Sigma \pm \frac{\partial^2 T}{\partial t_1 \partial t_1}\,\frac{\partial^2 T}{\partial t_2 \partial t_2}\cdots\frac{\partial^2 T}{\partial t_n \partial t_n},}

so ist

\weq{16}{H' = r^{-2}\,H.}

Wenn wir aber $H'$ nach (14) in den Variablen $\lambda_\nu$ bilden, und beachten, daß eine $n$-reihige Determinante, in der die Diagonalglieder $=0$ und alle anderen Glieder $=1$ sind, den Wert $(-1)^{n-1}(n-1)$ hat\footnote{Man kann diesen Satz leicht beweisen, wenn man die Determinante durch Zufügung einer $(n+1)$ten Zeile und einer $(n+1)$ten Spalte erweitert, bei der in der hinzugefügten Zeile nur \gesperrt{Einer}, in der Spalte mit Ausnahme des letzten Elementes Nullen stehen.}, so folgt:
%CONT>>
